% Generated draft: requires mathematical and source review.
\documentclass[11pt]{article}
\usepackage[margin=25mm]{geometry}
\usepackage{fontspec}
\setmainfont{Latin Modern Roman}
\usepackage{amsmath,amssymb,mathtools,unicode-math}
\setmathfont{Latin Modern Math}
\usepackage{enumitem,needspace,xurl,hyperref}
\setlength{\parindent}{0pt}
\setlength{\parskip}{5pt}
\setlength{\emergencystretch}{3em}
\newcommand{\N}{\mathbb{N}}
\newcommand{\Z}{\mathbb{Z}}
\newcommand{\Q}{\mathbb{Q}}
\newcommand{\R}{\mathbb{R}}
\newcommand{\C}{\mathbb{C}}
\newcommand{\F}{\mathbb{F}}
\newcommand{\A}{\mathbb{A}}
\newcommand{\PP}{\mathbb{P}}
\newcommand{\E}{\mathbb{E}}
\newcommand{\eps}{\varepsilon}
\newcommand{\dd}{\,\mathrm{d}}
\begin{document}
\section*{UnsolvedMath ID: 800880}
\textbf{Source problem code:} catalogue entry
\subsection{Definitions and mathematical statement}
Let \(\pi\colon X\to\mathbb{CP}^{2}\) be the blow-up of the complex projective plane at one point. Thus \(X\) is a del Pezzo surface with positive first Chern class. Fix a K\"ahler form \(\omega\) with \([\omega]=c_1(X)\), and a smooth anticanonical divisor \(D\subset X\), meaning a smooth effective divisor with \(\mathcal O_X(D)\simeq K_X^{-1}\). Write \([D]\) for the current of integration over \(D\). Set
\[
V:=\int_X\omega^2=\int_D\omega,
\qquad
\mu_D:=\frac{\omega\wedge[D]}{V}.
\]
The equality of the two integrals follows from \([D]=c_1(X)=[\omega]\) in cohomology. In particular, \(\mu_D\) is a probability measure characterized by
\[
\int_X f\,\dd\mu_D=\frac1V\int_D f|_D\,\omega|_D
\qquad(f\in C(X)).
\]
To make the notation explicit, adopt the convention \(dd^c=\frac{i}{\pi}\partial\bar\partial\). An \(\omega\)-plurisubharmonic function is an upper semicontinuous, locally integrable function \(u\colon X\to[-\infty,\infty)\), not identically \(-\infty\), which is locally the sum of a plurisubharmonic function and a smooth function and satisfies \(\omega+dd^c u\geq0\) as a current. Denote this class by \(\operatorname{PSH}(X,\omega)\). For bounded \(u\in\operatorname{PSH}(X,\omega)\), use the Bedford--Taylor product to define
\[
\operatorname{MA}_{\omega}(u):=\frac{(\omega+dd^c u)^2}{V}.
\]
Here \(\operatorname{DMA}(X,\omega)\) is interpreted as the global domain of definition under decreasing bounded approximation: it consists of those \(u\in\operatorname{PSH}(X,\omega)\) for which there exists a probability measure \(\nu_u\) such that, for every sequence \(u_j\in\operatorname{PSH}(X,\omega)\cap L^\infty(X)\) decreasing pointwise to \(u\),
\[
\operatorname{MA}_{\omega}(u_j)\rightharpoonup\nu_u.
\]
Set \(\operatorname{MA}_{\omega}(u)=\nu_u\) for such \(u\).

\textbf{Question (attributed in the source to Phong).} For every choice of \(\omega\) and \(D\) as above, does there exist \(\varphi_D\in\operatorname{DMA}(X,\omega)\) such that
\[
\operatorname{MA}_{\omega}(\varphi_D)=\mu_D?
\]
The supplied excerpt does not define \(V\), \(\operatorname{DMA}\), the normalization of \(\operatorname{MA}\), or the convention for \(dd^c\). The definitions above give an explicit standard interpretation requiring confirmation against the source's surrounding definitions. In particular, the operator here is defined by decreasing bounded approximation; it is not being interpreted as the non-pluripolar Monge--Amp\`ere product.
\subsection{Short English statement}
On the projective plane blown up at one point, fix a K\"ahler form representing the first Chern class and a smooth anticanonical divisor. The question asks whether normalized area measure on that divisor is the Monge--Amp\`ere measure of a potential in the global domain of definition. The operator and normalization conventions made explicit above require source review.
\subsection{Sources}
\begin{itemize}
\item Local UnsolvedMath registry, permanent ID \(800880\), exact source code \texttt{catalogue entry}, title: \emph{Open Problems in Pluripotential Theory --- Question 11}. The supplied registry labels this problem \texttt{open}; that status is not independently verified here.
\item Dinew, Guedj, Zeriahi, \emph{Open Problems in Pluripotential Theory} (2015), Question 11. Supplied source URL: \url{https://arxiv.org/abs/1511.00705}. The question itself carries the attribution \textnormal{[Phong]}; the supplied record does not specify a full name or a separate reference for that attribution.
\item Supplied provenance reports access on \texttt{2026-10-16} and extraction by \texttt{arxiv\_latex}. These are supplied provenance claims, not evidence of independent source consultation for this draft. The supplied upstream revision is \texttt{d371ee917ac9647fbb06cd6fc529deec6f3dda82}.
\end{itemize}
\end{document}
